\section{Conclusion}
This study addresses the post-optimization choice problem arising in large-scale maintenance scheduling. Rather than treating the best solver score as the final answer, the set of high-quality schedules is evaluated as a finite choice set under semantically meaningful criteria, and robust preference is assessed under attitudinal uncertainty.

The main methodological result is an exact equivalence, valid for a fixed finite choice set, that represents WOWA as a high-dimensional OWA on a common simplex. Its practical value does not lie in making WOWA computable for a given parameterization, since such point evaluations were already possible. Its value lies in making robust selection under WOWA analyzable in a shared linear weight space: alternative-specific cumulative-importance partitions are replaced by a common grid, indifference sets become hyperplanes, and the stability radius can be computed exactly by Linear Programming.

The application to the ROADEF/EURO 2020 schedules shows that the alternative favored by the competition score need not be the most defensible operational decision. Once technical priorities and attitudinal uncertainty are incorporated, preferred schedules may differ across strategic scenarios, and more conservative profiles tend to favor alternatives with stronger reliability properties. The framework not only provides an aggregate evaluation, but also an explicit measure of the stability of the selected schedule under admissible preference variations.

Future research should integrate the stability criterion into the search phase, refine the sustainability criteria with explicit geospatial data, and investigate how the unified weight-space representation induced by the theorem can support the transfer of elicitation methods from OWA-based MCDM models to WOWA and related rank-dependent aggregation settings.
